The study of complex physical systems commonly distinguishes between \textit{forward}
and \textit{inverse} problems. As a brief refresher, a forward problem consists of applying
governing physical laws—such as those of electromagnetism, mechanics, or physiology—to
a known set of inputs or system parameters in order to predict observable outcomes~\cite{mueller-2012}. 
Forward problems are typically well posed, admitting solutions that
are unique and stable with respect to perturbations in the input data~\cite{burger}.

Inverse problems address the complementary task: given a set of observations, the objective
is to reconstruct the causes, parameters, or internal states that may have produced those
measurements~\cite{mueller-2012}. While the term \textit{inverse problem} suggests a direct
mathematical reversal of the corresponding forward formulation, this relationship is not
guaranteed in general~\cite{burger}. In practice, inverse formulations often fail to satisfy
one or more of the classical criteria for well-posedness.

When modeling physical systems, these formulations are commonly interpreted in a
descriptive sense, where the emphasis lies on representing or explaining the behavior of
the system itself.

In a control-oriented setting, such as the one considered in the present work, the role of
forward and inverse problems is subtly different. The forward problem describes the
mapping between a set of accessible input parameters and the resulting system outputs,
thereby characterizing the system’s input–output behavior. The inverse problem, in turn,
seeks to determine which input configuration produces a specified output response. The
objective is not to recover the true physical origin of the observations, but to identify
admissible system inputs that reproduce or approximate a desired measurement pattern.

According to Hadamard, a mathematical problem is well-posed if a solution exists, is unique, and depends continuously on the input data~\cite{burger}. In system reconstruction contexts, the stability condition is often the most critical and the most frequently violated. When stability is lacking, even minute errors arising from measurement noise, discretization, or numerical round-off can be severely amplified during inversion, yielding reconstructions that are physically meaningless~\cite{burger}.

The transition from a continuous physical description to a computational model introduces discretization, reducing the governing differential or integral equations to a finite-dimensional numerical system. This process typically yields a linear algebraic representation of the form
\[
y = Ax,
\]
where $y$ denotes the vector of observations, $x$ represents the unknown system parameters or states, and $A$ is the system matrix (or forward operator) encoding the discretized physical model~\cite{mueller-2012}. The feasibility and stability of reconstruction or control-oriented tasks are therefore governed by the mathematical properties of this operator~\cite{mueller-2012, burger}.

In finite-dimensional settings, the operator $A$ arises from a mapping between an actuator space $X$ and an observation space $Y$, leading to a matrix $A \in \mathbb{R}^{m \times n}$, where $n$ denotes the number of parameters and $m$ the number of measurements~\cite{mueller-2012}. The relative dimensions of these spaces fundamentally determine the nature of the reconstruction problem. If $m = n$ and $A$ is nonsingular, the system admits a unique solution through the exact inverse $A^{-1}$~\cite{abdumuminova-2025}. Such ideal conditions, however, are rarely encountered in practical engineering applications.

Instead, two non-ideal regimes dominate:
\begin{enumerate}
    \item \textbf{Overdetermined Systems $m > n$:} There are more measurements than parameters. While this redundancy can be exploited to mitigate measurement noise, it generally results in inconsistent systems for which no exact solution $x$ satisfies $y = Ax$. Reconstruction in this regime is therefore formulated in a least-squares sense.

    \item \textbf{Underdetermined Systems $m < n$:} There are fewer measurements than  parameters, and the available data are insufficient to uniquely determine every degree of freedom of the system. Consequently, infinitely many solutions satisfy the measurement constraints, and additional criteria are required to select a physically meaningful solution.
\end{enumerate}

This work focuses on the overdetermined regime, in which the system matrix $A$ is non-square and therefore not invertible. In such cases, reconstruction must be reformulated in terms of a \textit{generalized solution} satisfying an auxiliary optimality criterion~\cite{mueller-2012}.

% \subsubsection{The Moore--Penrose pseudoinverse}
% As a robust generalization of the matrix inverse, the Moore--Penrose pseudoinverse $A^\dagger$ provides a unique solution for any matrix A, regardless of its dimensions or rank.
% The matrix $G \in \mathbf{R}^{n \times m}$ is said to be the Moore--Penrose pseudoinverse of $A$ if and only if it satisfies the following four algebraic conditions:
% \begin{enumerate}
%     \item $A\space G \space A = A$, maps the column vectors of $A$ back to themselves;
%     \item $G\space  A\space G = G$, establishing $G$ as a weak inverse of $A$;
%     \item $(A\space G)^* = A\space G$, implying that $AG$ is Hermitian;
%     \item $(G\space A)^* = G\space A$, implying that $GA$ is Hermitian.
% \end{enumerate}
% Here, $(\cdot)^*$ denotes the conjugate transpose. When considering real matrices, $(\cdot)^* = (\cdot)^T$, which implies the matrix is symmetric. While multiple generalized inverses may satisfy the first condition alone, the additional Hermitian constraints enforce orthogonality properties that guarantee the uniqueness of the Moore--Penrose pseudoinverse~\cite{penrose-1955}.

\subsubsection{Solutions in Overdetermined Systems}

When the number of observations exceeds the number of available actuators ($m > n$), the linear system
\[
y = Ax, \qquad A \in \mathbb{R}^{m \times n},
\]
is generally overdetermined and does not admit an exact solution for arbitrary measurement vectors $y$. In such cases, reconstruction must be interpreted in an approximate sense.

% The Moore--Penrose pseudoinverse provides a principled mechanism for defining an optimal solution by reformulating the problem as a least-squares minimization:
% \[
% x^\star = \arg\min_{x \in \mathbb{R}^n} \|Ax - y\|_2^2.
% \]
% The unique minimizer of this problem is given by
% \[
% x^\star = A^{\dagger} y.
% \]

% For full column-rank matrices, where $\operatorname{rank}(A) = n$, the pseudoinverse admits the explicit form
% \[
% A^{\dagger} = (A^{T}A)^{-1}A^{T},
% \]
% which corresponds to the classical normal-equations solution. Geometrically, this operation projects the observation vector $y$ orthogonally onto the column space of $A$, yielding the actuator configuration whose forward response best approximates the measurements in an energy-optimal sense~\cite{bindel-2012}.
A common approach is to define the solution as the parameter vector that minimizes the discrepancy between the predicted and observed outputs in a least-squares sense~\cite{mueller-2012, burger}. This leads to the optimization problem
\[
x^\star = \arg\min_{x \in \mathbb{R}^n} \|Ax - y\|_2^2.
\]

The minimizer of this problem can be expressed using the Moore--Penrose pseudoinverse of the system matrix, denoted $A^\dagger$, yielding
\[
x^\star = A^\dagger y.
\]
The Moore--Penrose pseudoinverse defines a generalized inverse applicable to arbitrary matrices, including non-square and rank-deficient cases. It is uniquely characterized by a set of algebraic conditions—known as the Moore--Penrose conditions—which enforce consistency with the original system and orthogonality properties of the associated projection operators, thereby guaranteeing uniqueness~\cite{penrose-1955}. In the present context, it serves as the operator that maps observations to the least-squares optimal reconstruction.

For matrices with full column rank, where $\operatorname{rank}(A) = n$, the pseudoinverse admits the explicit form
\[
A^{\dagger} = (A^{T}A)^{-1}A^{T},
\]
which corresponds to the classical normal-equations solution~\cite{bindel-2012}.


% In physical systems, this interpretation is particularly relevant: the pseudoinverse does not attempt to reproduce all observations exactly, but instead identifies the actuator inputs that minimize the global mismatch between predicted and measured outputs. The resulting residual reflects components of the observations that lie outside the span of the system’s controllable subspace.

\subsubsection{Regularization and Numerical Stability}

While $A^{\dagger}$ provides the optimal least-squares solution, it may be unstable in the presence of noise or modeling errors for ill-conditioned systems. A system is ill-conditioned if its \textbf{condition number} $\kappa(A) = \sigma_{\max}/\sigma_{\min}$ is large, where $\sigma_{\max}$  and $\sigma_{\min}$ denote the largest and smallest singular values of $A$.

In such cases, small perturbations in the observation vector $y$ can be strongly amplified in the reconstructed solution, since $\|A^{\dagger}\|_2 = 1/\sigma_{\min}(A)$~\cite{golub-2013, mueller-2012}.

When a system is ill-conditioned, the naive application of the exact pseudoinverse can introduce error. Regularization must be used to stabilize the solution by introducing a bias toward physically plausible outcomes~\cite{mueller-2012}. In a broader system-design context, such a bias may also be interpreted as a deliberate modeling or control choice, reflecting desired properties of the solution and allowing a controlled trade-off between output accuracy and feasibility.


Tikhonov regularization (also known as $L_2$ regularization) transforms the ill-posed least-squares problem into a well-posed penalized problem:


$$\min_{x} \left\{ \|Ax - y\|_2^2 + \lambda^2 \|L x\|_2^2 \right\}$$

where $\lambda$ is the regularization parameter and $L$ is the regularization operator.
The regularized solution $x_{\lambda}$ is:


$$x_{\lambda} = (A^T A + \lambda^2 L^T L)^{-1}A^T y$$


The additional term $\lambda^2 L^TL$ increases the smallest singular values of the system, thereby reducing noise amplification and improving numerical stability~\cite{golub-2013, mueller-2012, vara-kumar-2022}.

The most common choice is $L = I$, the identity matrix, which penalizes all components of the solution vector $x$ uniformly, favoring solutions with small overall energy. More generally, however, $L$ may be chosen to encode prior knowledge or design preferences about the system.